% Source: https://artofproblemsolving.com/community/c3h3131564_2023_emcc_team_round__exeter_math_club_competition
% Stable identity: aops-topic-3131564-q04
The first three numbers have sum $9$, while $b+c=8$ from the median, hence $a=1$. The overall range then gives $f=7$. Since the unique mode of the first five entries is $5$, at least two of those entries must equal $5$. The ordered possibilities for $(b,c)$ with $b+c=8$ are $(1,7),(2,6),(3,5),(4,4)$. In the first two cases every later entry is greater than $5$, so they cannot work. In the last case, two copies of $4$ already occur, and at most two copies of $5$ could occur, giving a tie rather than unique mode $5$. Thus $(b,c)=(3,5)$. We must have $d=5$, while $e$ may be $5,6,$ or $7$; in each case $5$ is indeed the unique mode of the first five entries. The product $1\cdot3\cdot5\cdot5\cdot e\cdot7$ is greatest at $e=7$. The tuple $(1,3,5,5,7,7)$ attains $\boxed{3675}$.
